\documentclass[12pt]{article}
\usepackage[a4paper,
	left=2cm,
	right=2.5cm,
	top=3cm,
	bottom=3cm]{geometry}
\usepackage[a-1b]{pdfx}
\usepackage{amsmath}
\usepackage{amssymb}
\usepackage{amsthm}
\usepackage{mathtools}
\begin{document}

The rigorous analysis of a many-body quantum systems poses numerous hard challenges, as the
complexity of the wavefunction grows very rapidly in the number of degrees of freedom it
depends on. Extracting results about the energy spectrum or the structure of the ground state
can prove quite difficult and requires understanding the combined effects of particle
statistics, non-trivial correlations and collective behaviour.
In many cases, a natural way to account for these mechanisms is to identify regimes in which
the full, interacting theory can be replaced -- at least approximately -- by a simpler,
effective theory. The idea is that not all microscopic features of a system contribute equally
to the macroscopic behavior, and one may look for suitable asymptotic regimes in which
the collective effect of many particles can be described by effective quantities that can
be treated directly.

An example of such regimes, and the one which we will explore in this thesis, is that of the
\textit{mean field}. In this regime, the number of particles in the system is
very large, and the interaction between them is rescaled such that the total effect on the
energy of a single particle remains of order one. Even in the mean field limit, however, the
physical properties of a system may still be hard to understand, and this motivates the
perturbative approach. One starts from a sufficiently well understood problem, such as a
non-interacting or mean field regime, and introduces corrections to study the effects of the
interactions on the true interacting ground state.

In the present thesis we analyze the ground state energy of an interacting
Fermi gas in the mean field limit. This limit is realized by choosing the physical parameters
in the Hamiltonian so that they scale with the number of particles $N$. In particular, the
kinetic energy of the particles is defined in terms of a semi-classical parameter $\hbar$,
chosen to be of order $N^{-1/3}$, while the two-particle interaction is defined in terms of
the coupling paramter $\lambda$, chosen to be of order $N^{-1}$.
We follow a perturbative approach to compute the free energy of the system up to second order,
show that the ground state energy reproduces the scaling in $N$ present in the literature
up to first order, and discuss the first steps taken towards second order perturbation theory.

\end{document}
